\scanpage{92}
\mspage{45}
complet de $X'$ se déduit de celui de $X$, en prenant la $=$ application genre, et $I'\longrightarrow S$ induit par $I\longrightarrow S$.

En fait, quand on a normalisé $\widehat X$ en les pts de $\unsure{A|A^{a}}\setminus A'$ pour obtenir $\widehat X'$, on a sur $\widehat X'$, en plus des pts marqués indexés par $I$ et des pts \deleted{quadratiques}\inserted{doubles} indexés par $A'$, un troisième ens. remarquable de « pts marqués », savoir l'image inverse \inserted{$\unsure{\widetilde J|\widehat J}$} de $J=A\setminus A'$, \deleted{E} de cardinal égal à $2\operatorname{Card}(A\setminus A')$. En termes de la courbe normale propre $Y$ \deleted{qui} contenant $\widetilde A\amalg I$, qui définit $\widehat X$ comme $Y/\sigma_{\widetilde A}$ et $X$ comme $(Y\setminus I)/\sigma_{\widetilde A}$, $\unsure{\widehat X'|X'}$ est défini comme $Y/\sigma_{\widetilde A'}$, ce qui nous donne un ens. de pts marqués lisses \deleted{$\unsure{A}$} $\widetilde J\amalg I$ sur $\widehat X'$. Ceci étant, on peut « oublier » $\widetilde J$ et $I\setminus I'$ pour ne retenir que l'ens. de pts marqués
\margin{91}
